\documentclass[11pt,a4paper]{article}
% Common preamble for the hint sheets (\input from each file)
\usepackage[utf8]{inputenc}
\usepackage[T1]{fontenc}
\usepackage{lmodern}
\usepackage{amsmath,amssymb,amsthm,mathtools}
\usepackage[a4paper,margin=2.3cm]{geometry}
\usepackage{enumitem}
\usepackage{xcolor}
\usepackage[hidelinks]{hyperref}
\renewcommand{\leq}{\leqslant}
\renewcommand{\le}{\leqslant}
\renewcommand{\geq}{\geqslant}
\renewcommand{\ge}{\geqslant}
\setlength{\parindent}{0pt}
\setlength{\parskip}{0.35em}
% \hintstitle{Sheet title}{Class work / Homework}
\newcommand{\hintstitle}[2]{\begin{center}{\LARGE\bfseries Hints: #1\par}\vspace{0.5em}{\large #2\par}\end{center}
\noindent\textit{The hints to each problem are cumulative. Read one, try again, and open the next only if you are still stuck.}\par\medskip\hrule\medskip}
\newcommand{\problem}[1]{\par\bigskip{\large\bfseries Problem #1}\par\nopagebreak\smallskip}
\newcommand{\subpart}[1]{\par\smallskip\textbf{(#1)}\par\nopagebreak}
\newcommand{\hint}[1]{\par\smallskip\noindent\textcolor{black!70}{\textbf{Hint #1.}}\enspace}
\newcommand{\note}[1]{\par\smallskip\noindent\textit{Note.} #1\par}
\providecommand{\lt}{<}
\providecommand{\gt}{>}
\providecommand{\tightlist}{\setlength{\itemsep}{0pt}\setlength{\parskip}{0pt}}
\emergencystretch=3em

\begin{document}
\hintstitle{Polynomials -- 7. Interpolation and finite differences}{Homework}

\problem{1}
\hint{1} Write $f$ in Newton's form $f(x)=\sum_k\Delta^kf(0)\binom xk$ (both sides have degree $\le\deg f$ and agree at $0,1,\dots,\deg f$). If $f$ is integer-valued, each $\Delta^kf(0)$ is an integer combination of the values $f(0),\dots,f(k)$.
\hint{2} Examples: $\binom xn=\frac{x(x-1)\cdots(x-n+1)}{n!}$ is integer-valued of degree $n$ with leading coefficient $1/n!\notin\mathbb{Z}$ for every $n\ge2$; e.g.\ $\binom x2=\frac{x^2}2-\frac x2$. (The constant term $f(0)$ is always an integer.)

\problem{2}
\hint{1} Put $m=\lfloor n/k\rfloor$ and $r=n-km\in\{0,\dots,k-1\}$. Find a polynomial identity in $n$ and $m$ that holds because $r$ is one of $0,\dots,k-1$.
\hint{2} $\prod_{j=0}^{k-1}(n-km-j)=0$, since one factor is $r-r$. Expanded in powers of $m$, this is $(-k)^km^k+\sum_{i<k}Q_i(n)m^i=0$ with polynomials $Q_i$; divide by $-(-k)^k$.

\problem{3}
\hint{1} Write $g_n=\sum_jc_{n,j}x^j$. The rule says $c_{n+1,j}=c_{n,j-1}+c_{n,j+1}$, except that nothing is carried to exponent $-1$. Interpret $c_{n,j}$ as a number of lattice paths.
\hint{2} $c_{n,j}$ counts walks of $n$ steps $\pm1$ from $0$ to $j$ that never go below $0$. So $g_n(0)$ is the number of Dyck paths of length $n$: $0$ for odd $n$ and the Catalan number $C_{n/2}$ for even $n$.
\hint{3} Answer: $g_{15}(0)=0$ and $g_{16}(0)=C_8=\frac1{9}\binom{16}8=1430$.

\problem{4}
\hint{1} All differences of $2^k$ equal $2^k$ again, so $\Delta^j(2^k)\big|_{k=0}=1$. Write $f$ in Newton's form.
\hint{2} $f(x)=\sum_{j=0}^n\binom xj$. Then $f(n+1)=\sum_{j\le n}\binom{n+1}j=2^{n+1}-1$ and $f(n+2)=\sum_{j\le n}\binom{n+2}j=2^{n+2}-n-3$. In terms of differences: $f$ agrees with $2^k$ except that its $(n+1)$-st difference is $0$ instead of $1$.

\problem{5}
\hint{1} Bound $\max_{[0,1]}|P|$ from below by an integral that is a positive rational number with controlled denominator. Try $\int_0^1P(x)^{2m}\,dx$ for large $m$.
\hint{2} $\int_0^1P^{2m}$ is positive and a sum of terms (integer)$/(j+1)$ with $j\le2mn$, hence at least $1/\operatorname{lcm}(1,\dots,2mn+1)=e^{-\psi(2mn+1)}$, where $\psi$ is Chebyshev's function. So $\max_{[0,1]}|P|\ge\exp\bigl(-\psi(2mn+1)/(2m)\bigr)$ for every $m$.
\hint{3} Let $m\to\infty$: since $\psi(N)/N\to1$ (prime number theorem; in fact $\liminf\psi(N)/N\le1$ suffices, which follows from Mertens' estimate $\sum_{j\le N}\Lambda(j)/j=\ln N+O(1)$), the right-hand side gets arbitrarily close to $e^{-n}$ along a subsequence, hence $\max|P|\ge e^{-n}$.

\problem{6}
\hint{1} Forward difference formula: $f(x)=\sum_k\Delta^kf(0)\binom xk$, proved by checking both sides agree at $x=0,\dots,\deg f$ (or by $E^x=(I+\Delta)^x$ on integers). Combine with $\sum_{k=0}^{m}\binom kj=\binom{m+1}{j+1}$.
\hint{2} $k^3=6\binom k3+6\binom k2+\binom k1$ (differences of $k^3$ at $0$: $0,1,6,6$), so $\sum_{k=1}^mk^3=6\binom{m+1}4+6\binom{m+1}3+\binom{m+1}2=\bigl(\frac{m(m+1)}2\bigr)^2$. Also $k(k+1)(k+2)=6\binom{k+2}3$, so the second sum is $6\binom{m+3}4=\frac{m(m+1)(m+2)(m+3)}4$.

\problem{7}
\hint{1} $P(x)-7=(x-a_1)(x-a_2)(x-a_3)(x-a_4)Q(x)$ with $Q\in\mathbb{Z}[x]$.
\hint{2} If $P(k)=14$, then $k-a_1,\dots,k-a_4$ are four distinct integers whose product divides $7$. The only divisors of $7$ are $\pm1,\pm7$, and the product of four distinct ones is $49$, which does not divide $7$.
\hint{3} General statement: if $P\in\mathbb{Z}[x]$ takes the value $c$ at four (or more) distinct integers, then $P(k)-c$ is never $\pm1$ or $\pm$ a prime (more generally, $P(k)-c$ is divisible by a product of four distinct nonzero integers).

\problem{8}
\note{As stated (nodes arbitrary in $[0,1]$) the bound is false: for the nodes $0,\varepsilon$ and values $0,1$ the interpolant $x/\varepsilon$ has an unbounded coefficient. The constant must also depend on the nodes, e.g.\ through the minimal distance $\delta$ between them; the hints follow this version.}
\hint{1} Vandermonde: $\det(x_i^{j})$ is a polynomial in $x_n$ of degree $n$ (say $n+1$ nodes $x_0,\dots,x_n$) vanishing at $x_0,\dots,x_{n-1}$, so it equals the smaller Vandermonde determinant times $\prod_{i<n}(x_n-x_i)$; induct.
\hint{2} The coefficient vector is $V^{-1}y$. By Cramer's rule each coefficient is $\det V_k/\det V$, where $V_k$ has entries in $[0,1]$ except one column equal to $y$, so $|\det V_k|\le(n+1)!\max|y_i|$, while $|\det V|=\prod_{i<j}|x_j-x_i|\ge\delta^{n(n+1)/2}$.

\problem{9}
\hint{1} By Lagrange, $|P(x)|\le\sum_{k=0}^n|L_k(x)|$ with $L_k(x)=\prod_{j\ne k}\frac{x-j}{k-j}$ and $\prod_{j\ne k}|k-j|=k!\,(n-k)!$. It suffices to show $\prod_{j\ne k}|x-j|\le n!$ for $x\in[0,n]$; then the sum is at most $\sum_k\binom nk=2^n$.
\hint{2} Let $x\in[m,m+1]$, $0\le m\le n-1$, and by symmetry $k\le m$. Bound $|x-j|\le j-m$ for $j>m$ and $|x-j|\le m+1-j$ for $j\le m$: the product is at most $(n-m)!\,(m+1)!/(m+1-k)\le(m+1)!\,(n-m)!$, and this is $\le n!$ because $\binom nm\ge m+1$.
\hint{3} Sharpness: take $P(k)=(-1)^k$. Then $\Delta^nP=(-2)^n$, so the leading coefficient is $(-2)^n/n!$, and by Chebyshev a polynomial with leading coefficient $a$ has $\sup_{[0,n]}|P|\ge2|a|(n/4)^n$. This gives $\sup|P|\ge2(n/2)^n/n!$, which grows like $(e/2)^n/\sqrt n$.

\problem{10}
\hint{1} Define $P[x_0,\dots,x_k]$ as the leading coefficient (of $x^k$) of the interpolating polynomial $L_k$ through $(x_i,P(x_i))$, $i\le k$. Reading off the $x^k$ coefficient of Lagrange's formula gives the stated symmetric formula; check it satisfies the recursion $P[x_0,\dots,x_k]=\frac{P[x_1,\dots,x_k]-P[x_0,\dots,x_{k-1}]}{x_k-x_0}$.
\hint{2} Newton: $L_k-L_{k-1}$ has degree $\le k$ and vanishes at $x_0,\dots,x_{k-1}$, so it equals $P[x_0,\dots,x_k]\prod_{j<k}(x-x_j)$. Summing, $L_n(x)=\sum_{k\le n}P[x_0,\dots,x_k]\prod_{j<k}(x-x_j)$.
\hint{3} If $\deg P=n$, the interpolant through $n+1$ nodes is $P$ itself, so its order-$n$ divided difference is the leading coefficient of $P$.

\end{document}
